% New section: kernel normal form and the exponential defect.

\section{Kernel normal form and the exponential defect}
\label{sec:kernel}

Throughout this section we work at an ordinary primitive canonical seed
with the shared local neighborhood and based representatives of
Definition~\ref{def:shared-neighborhood}, and we write
\[
 a=i^{p}\colon\Omega_s\hookrightarrow\Omega_t,
 \qquad
 N_{\mathrm{fill}}:=\ker\bigl[a_*\colon\pi_1(\Omega_s)\to\pi_1(\Omega_t)\bigr],
 \qquad
 K=a_*^{-1}(H_t).
\]
We call $a_*$ the \emph{hole-filling homomorphism}: it is induced by the
inclusion that fills back into the target every component of
$F^{-1}B\setminus B$ removed from the source.

\subsection{The sheet fiber is the hole-filling kernel}

\begin{theorem}[Kernel normal form]\label{thm:kernel-normal-form}
At an ordinary primitive canonical seed, the restriction
$a_*|_{H_s}\colon H_s\to H_t$ is an isomorphism.  Consequently
\[
 K=N_{\mathrm{fill}}\rtimes H_s,
\]
and the map
\[
 N_{\mathrm{fill}}\longrightarrow K/H_s,
 \qquad
 n\longmapsto nH_s,
\]
is a canonical pointed bijection.
\end{theorem}

\begin{proof}
By Lemma~\ref{lem:local-pair}, $D_A\setminus A$ is an annulus; let $c$
be its positively oriented core generator, so
$\pi_1(D_A\setminus A)\cong\mathbb Z$ is generated by $[c]$.  The image
of $[c]$ in $\pi_1(\Omega_s)$, and likewise in $\pi_1(\Omega_t)$, has
winding number one about every point of $A$; since the winding number
about a fixed point of $A$ is a homomorphism to $\mathbb Z$, the image
of $[c]^{n}$ is nontrivial for every $n\neq0$.  Hence both induced
maps
\[
 \pi_1(D_A\setminus A)\cong\mathbb Z\longrightarrow H_s,
 \qquad
 \pi_1(D_A\setminus A)\cong\mathbb Z\longrightarrow H_t
\]
are isomorphisms onto their images.  Because $a$ restricts to the
identity on the common physical neighborhood $D_A\setminus A$, the
triangle formed by these two maps and $a_*|_{H_s}$ commutes, and
$a_*|_{H_s}\colon H_s\to H_t$ is an isomorphism.

Let $k\in K$.  Since $a_*(k)\in H_t$, there is a unique $h\in H_s$ with
$a_*(h)=a_*(k)$; then $n:=kh^{-1}\in\ker a_*=N_{\mathrm{fill}}$ and
$k=nh$.  Injectivity of $a_*|_{H_s}$ gives
$N_{\mathrm{fill}}\cap H_s=\{1\}$, and $N_{\mathrm{fill}}$, being the
kernel of $a_*$ on $\pi_1(\Omega_s)$, is normal in $K$.  Hence
$K=N_{\mathrm{fill}}\rtimes H_s$.  Finally, $k=nh$ gives
$kH_s=nH_s$, and if $nH_s=n'H_s$ with
$n,n'\in N_{\mathrm{fill}}$ then
$n^{-1}n'\in N_{\mathrm{fill}}\cap H_s=\{1\}$; so
$n\mapsto nH_s$ is a bijection onto $K/H_s$, pointed because
$1\mapsto H_s$.
\end{proof}

\begin{remark}\label{rem:kernel-interpretation}
The pointed coset set $K/H_s$ of
Theorem~\ref{thm:sheet-resolution} is not in general a group.
Theorem~\ref{thm:kernel-normal-form} shows that at the ordinary
canonical seed, with the shared based representatives of
Definition~\ref{def:shared-neighborhood}, it nevertheless carries a
canonical group model: the kernel of the hole-filling homomorphism.
In one sentence, \emph{canonical localization forgets exactly the
kernel of the source-to-target hole-filling map}.  The pointed-coset
discipline of Sections~\ref{sec:resolution}--\ref{sec:enlargement}
remains the correct general-position language; the group model is a
feature of the canonical seed with the shared neighborhood, not of the
abstract setting.
\end{remark}

\subsection{The exact excess count}

\begin{theorem}[Exact excess count]\label{thm:excess-count}
Let $f(z)=z^{2}+c$ be primitively renormalizable of period $p>1$ with
ordinary return $F=f^{p}\colon U^{p}\to U^{0}$ and central continuum
$A=K_p$.  Then
\[
 F^{-1}(A)=A\sqcup X_1\sqcup\cdots\sqcup X_{2^{p}-2},
\]
where each $X_j$ is a full continuum disjoint from $B$, and
$F|_{X_j}\colon X_j\to A$ is a homeomorphism.
\end{theorem}

\begin{proof}
Differentiating the composition,
\[
 F'(z)=\prod_{r=0}^{p-1}f'\bigl(f^{r}(z)\bigr)
      =2^{p}\prod_{r=0}^{p-1}f^{r}(z),
\]
so $z$ is a critical point of $F$ exactly when $f^{r}(z)=0$ for some
$0\le r<p$, with critical value
$F(z)=f^{p-r}\bigl(f^{r}(z)\bigr)=f^{p-r}(0)$.  Since $0\in A$ and
$f^{j}(A)\subset K_{j}$, we get $f^{p-r}(0)\in K_{p-r}$.  For
$1\le r<p$ the index $p-r$ lies in $\{1,\dots,p-1\}$, so by primitive
disjointness the critical value lies in a member of $B\setminus A$ and
in particular not in $A$.  For $r=0$ the critical point is $z=0$
itself, with critical value $f^{p}(0)\in A$.  Hence \emph{the only
critical point of $F$ whose critical value lies in $A$ is $0$}.

\emph{Step 1: isolate $A$ in the source.}  By
Corollary~\ref{cor:finite-pullback} the preimage $F^{-1}(A)$ has
finitely many components, each a full continuum.  Choose a Jordan
neighborhood $N_A\Subset U^{p}$ of $A$ whose closure misses every
component of $F^{-1}(A)$ other than $A$, with
$\partial N_A\cap F^{-1}(A)=\varnothing$.  Then $F(\partial N_A)$ is
compact and disjoint from $A$.

\emph{Step 2: choose the target neighborhood.}  Choose a Jordan
neighborhood $V\supset A$ with $\overline V\subset U^{0}$, small enough
that $\overline V$ is disjoint from $B\setminus A$, from every
noncentral critical value $f^{p-r}(0)$ ($1\le r<p$), and from
$F(\partial N_A)$.

\emph{Step 3: the central component.}  Let $V_0$ be the component of
$F^{-1}(V)$ containing $A$ (hence containing $0$).  Since $V_0$ is
connected, meets $N_A$, and satisfies
$F(V_0)\subset V$, while $V\cap F(\partial N_A)=\varnothing$, the set
$V_0$ cannot meet $\partial N_A$; hence $V_0\Subset N_A$.  Every
critical point of $F$ lying in $F^{-1}(V)$ has critical value in $V$,
hence equals $0$; so $V_0$ is the unique component containing a
critical point, and $F|_{V_0}\colon V_0\to V$ is proper of
degree~$2$.  Because $N_A$ contains no component of $F^{-1}(A)$ other
than $A$,
\[
 F^{-1}(A)\cap V_0\subset F^{-1}(A)\cap N_A=A ,
 \qquad\text{so}\qquad
 F^{-1}(A)\cap V_0=A .
\]

\emph{Step 4: the univalent components.}  Every other component $V_j$
of $F^{-1}(V)$ contains no critical point, so it carries a proper
unbranched covering of the simply connected domain $V$, hence a
biholomorphism.  The degrees over $V$ sum to $\deg F=2^{p}$, so there
are exactly $2^{p}-2$ biholomorphic components
$V_1,\dots,V_{2^{p}-2}$.

\emph{Step 5: the excess continua.}  For $j\ge1$ set
$X_j=(F|_{V_j})^{-1}(A)$: a homeomorphic copy of $A$, compact,
connected, and a component of $F^{-1}(A)$.

\emph{Step 6: fullness.}  Each $X_j$ is a component of the $p$-fold
pullback of the full continuum $A$, so it is full by
Corollary~\ref{cor:finite-pullback}.

\emph{Step 7: disjointness from $B$.}  This is the argument of
Proposition~\ref{prop:excess}: an intersection $X_j\cap K_l$ would map
into $A\cap K_l=\varnothing$ for $l\neq0$ by primitivity, while
$X_j\cap A=\varnothing$ because distinct components of $F^{-1}(A)$ are
disjoint.
\end{proof}

\begin{corollary}[Exponential homological defect]\label{cor:exponential-defect}
By Corollary~\ref{cor:finite-pullback} the set $F^{-1}(B)$ has
finitely many components, all full continua, among them $B$'s own
components and the excess continua $X_j$.  Apply
Corollary~\ref{cor:isolating-disks} to this \emph{entire} finite
disjoint family and retain the disks belonging to the $X_j$: this
yields pairwise disjoint isolating Jordan disks $D_j\supset X_j$ with
\[
 D_j\cap\bigl(F^{-1}(B)\setminus X_j\bigr)=\varnothing ,
\]
so in particular $D_j\cap B=\varnothing$ and
$\partial D_j\subset\Omega_s$.  Choose based paths
$\eta_j\subset\Omega_s$, peripheral loops
$\gamma_j=\eta_j*\partial D_j*\eta_j^{-1}$, and points $x_j\in X_j$.
Then each $\gamma_j$ represents a class of $N_{\mathrm{fill}}=\ker a_*$,
and the winding vector
\[
 W=\bigl(\operatorname{Ind}_{x_1},\dots,
         \operatorname{Ind}_{x_{2^{p}-2}}\bigr)\colon
 N_{\mathrm{fill}}\longrightarrow\mathbb Z^{\,2^{p}-2}
\]
is a surjective homomorphism with $W([\gamma_j])=e_j$.  Consequently
\[
 \operatorname{rank}H_1\bigl(N_{\mathrm{fill}};\mathbb Z\bigr)\;\ge\;2^{p}-2 ,
\]
and, via Theorem~\ref{thm:kernel-normal-form}, the canonical sheet
fiber $K/H_s$ carries at least $2^{p}-2$ independent peripheral winding
coordinates.
\end{corollary}

\begin{proof}
Each $x_j$ lies in $F^{-1}B$, so $x_j\notin\Omega_s$ and
$\operatorname{Ind}_{x_j}$ is a well-defined homomorphism
$\pi_1(\Omega_s)\to\mathbb Z$.  Each $D_j$ is disjoint from $B$, hence
contained in $\Omega_t$, and $\gamma_j$ contracts across $D_j$ there,
so $a_*[\gamma_j]=1$ and $[\gamma_j]\in N_{\mathrm{fill}}$.  Pairwise
disjointness of the disks gives
$\operatorname{Ind}_{x_i}([\gamma_j])=\delta_{ij}$, i.e.\
$W([\gamma_j])=e_j$; hence $W$ is surjective.  A surjection of
$N_{\mathrm{fill}}$ onto the free abelian group of rank $2^{p}-2$
factors through $H_1(N_{\mathrm{fill}};\mathbb Z)$, giving the rank
bound.
\end{proof}

\begin{remark}\label{rem:sharper-bookkeeping}
Corollary~\ref{cor:exponential-defect} sharpens
Corollary~\ref{cor:infinite-defect} from ``the sheet defect is
infinite'' to ``the sheet defect has at least $2^{p}-2$ independent
winding coordinates.''  Any finite bookkeeping of the canonical
localization must therefore lose not one but exponentially many
independent peripheral directions, at every ordinary primitive seed.
\end{remark}

\subsection{The unicritical analogue}

\begin{lemma}[Degree-$d$ pullback dichotomy]\label{lem:degree-d-pullback}
Let $g\colon U'\to U$ be a proper holomorphic map of degree $d$ between Jordan disks whose only critical point $\omega$ has local degree $d$, and let $K\subset U$ be a full continuum.  If $g(\omega)\notin K$, then $g^{-1}(K)$ has exactly $d$ components, each a full continuum mapped homeomorphically onto $K$.  If $g(\omega)\in K$, then $g^{-1}(K)$ is a single full continuum, mapped properly with degree $d$.
\end{lemma}
\begin{proof}
Choose nested Jordan neighborhoods $W_1\Supset W_2\Supset\cdots$ of $K$ with $\bigcap_k W_k=K$, chosen to avoid $g(\omega)$ when $g(\omega)\notin K$.  \emph{Unbranched case:} over each simply connected $W_k$ the restriction of $g$ is a proper unbranched covering, hence splits into $d$ biholomorphic inverse branches; the nested intersections of the branch images are $d$ disjoint full continua mapped homeomorphically onto $K$ (fullness passes to nested intersections of Jordan domains).  \emph{Branched case:} since the local model of $g$ at $\omega$ is $z\mapsto z^{d}$, the fiber $g^{-1}(g(\omega))$ is the single point $\omega$; as every component of $g^{-1}(W_k)$ maps onto $W_k$ by properness and $g(\omega)\in W_k$, every component contains $\omega$, so $g^{-1}(W_k)$ is connected and maps with full degree $d$.  By Riemann--Hurwitz with one ramification point of multiplicity $d-1$, its Euler characteristic is $d\cdot1-(d-1)=1$, so each $g^{-1}(W_k)$ is a simply connected Jordan domain; the nested intersection $g^{-1}(K)=\bigcap_k g^{-1}(W_k)$ is a single full continuum mapped properly with degree $d$.
\end{proof}

\begin{theorem}[Unicritical analogue]\label{thm:unicritical}
Let $f(z)=z^{d}+c$ with $d\ge2$ be primitively renormalizable of period
$p>1$, with ordinary return $F=f^{p}$ and central continuum $A$.  Then
\[
 F^{-1}(A)=A\sqcup X_1\sqcup\cdots\sqcup X_{d^{p}-d},
\]
with every $X_j$ a full continuum disjoint from $B$ and
$F|_{X_j}\colon X_j\to A$ a homeomorphism, and the corresponding
winding vector is a surjection
$N_{\mathrm{fill}}\twoheadrightarrow\mathbb Z^{\,d^{p}-d}$.
\end{theorem}

\begin{proof}
Now $F'(z)=d^{p}\prod_{r=0}^{p-1}\bigl(f^{r}(z)\bigr)^{d-1}$, so the
critical points of $F$ are again the solutions of $f^{r}(z)=0$,
$0\le r<p$, with critical values $f^{p-r}(0)$; exactly as in
Theorem~\ref{thm:excess-count}, only $z=0$ has critical value in $A$,
the central component of $F^{-1}(V)$ has degree $d$, and the degree
budget leaves exactly $d^{p}-d$ unbranched degree-one components.  The
isolation of $A$ by $N_A$, the choice of $V$ avoiding $F(\partial N_A)$
and the noncentral critical values, the fullness of the $X_j$ via the
degree-$d$ analogue of Corollary~\ref{cor:finite-pullback}
(Lemma~\ref{lem:degree-d-pullback} replacing the quadratic dichotomy;
the finite-full-pullback induction then applies verbatim), the disjointness
from $B$, and the disk isolation from the full finite family of
components of $F^{-1}(B)$ all repeat verbatim.  We phrase this as the unicritical
analogue of the forced-collision setting: the localization construction
of Definition~\ref{def:localization} is stated for general degree, but
the term \emph{canonical renormalization} is used by the source in the
quadratic setting, and we attribute no degree-$d$ canonical
renormalization definition to it.
\end{proof}

\subsection{Future work}

\begin{remark}[Toward the full kernel]\label{rem:free-kernel}
Let $r$ be the number of source-only components of
$F^{-1}(B)\setminus B$, so that $r\ge 2^{p}-2$ by
Theorem~\ref{thm:excess-count}.  A finite-connectivity argument with
compatible peripheral bases should identify
$\pi_1(\Omega_s)\cong F_{p+r}$ and $\pi_1(\Omega_t)\cong F_{p}$ (free
groups of the indicated ranks) and identify $a_*$ with the free
retraction killing the $r$ extra peripheral generators.  This would
make $N_{\mathrm{fill}}$ a free group of countably infinite rank and
sharpen the rank bound of Corollary~\ref{cor:exponential-defect} to
\[
 \operatorname{rank}H_1\bigl(N_{\mathrm{fill}};\mathbb Z\bigr)
 =\aleph_0 ;
\]
the winding vector $W$ would then be the canonical finite quotient
attached to the $2^{p}-2$ excess preimages of $A$.  We do not claim
this here: the finite-connectivity and peripheral-basis lemma is not
yet written.
\end{remark}
